\documentclass[11pt]{article}
\usepackage{amsmath, amssymb, amsthm}
\usepackage[margin=1in]{geometry}
\usepackage{algorithm}
\usepackage{algpseudocode}

\newtheorem{theorem}{Theorem}
\newtheorem{lemma}[theorem]{Lemma}

\title{Problem 411: Uphill Paths}
\author{}
\date{}

\begin{document}
\maketitle

\section{Problem Statement}

Let $n$ be a positive integer. Stations are placed at coordinates
\[
  (x, y) = (2^i \bmod n, \; 3^i \bmod n) \quad \text{for } 0 \le i \le 2n.
\]
Stations at identical coordinates are identified. An \emph{uphill path} from $(0,0)$ to $(n,n)$ is one where both coordinates never decrease. Let $S(n)$ denote the maximum number of stations such a path can pass through.

\textbf{Find:} $\displaystyle\sum_{k=1}^{30} S(k^5)$.

\section{Mathematical Foundation}

\begin{theorem}[Reduction to LNDS]
Let $\mathcal{P} = \{(x_1,y_1),\dots,(x_M,y_M)\}$ be the set of distinct stations, sorted lexicographically by $(x,y)$. Then $S(n)$ equals the length of the longest non-decreasing subsequence (LNDS) of $(y_1,\dots,y_M)$.
\end{theorem}

\begin{proof}
An uphill path visits stations in order of non-decreasing $x$-coordinate. Among all subsequences of stations with non-decreasing $x$, the maximum-length subsequence that also has non-decreasing $y$ is precisely the LNDS of the $y$-values when stations are sorted by $(x,y)$ lexicographically. The lexicographic sort with $y$ ascending as tiebreaker ensures that when two stations share the same $x$, both may appear on a non-decreasing path. Formally, a subsequence $(x_{i_1},y_{i_1}),\dots,(x_{i_\ell},y_{i_\ell})$ with $i_1 < \cdots < i_\ell$ satisfies $x_{i_1} \le \cdots \le x_{i_\ell}$ (by sort order) and $y_{i_1} \le \cdots \le y_{i_\ell}$ (by the LNDS property), which is exactly the uphill condition. Conversely, every uphill subsequence induces a non-decreasing subsequence of $y$-values in the sorted order.
\end{proof}

\begin{lemma}[Periodicity]
The sequence $2^i \bmod n$ has period $\mathrm{ord}_n(2) \mid \varphi(n)$, and similarly for $3^i \bmod n$. The joint station sequence has period $\mathrm{lcm}(\mathrm{ord}_n(2), \mathrm{ord}_n(3))$.
\end{lemma}

\begin{proof}
By Euler's theorem, $a^{\varphi(n)} \equiv 1 \pmod{n}$ when $\gcd(a,n)=1$, so $\mathrm{ord}_n(a) \mid \varphi(n)$. The sequence $a^i \bmod n$ is periodic with this order. Since $2n + 1 > \varphi(n)$ for all $n \ge 1$, all distinct stations appear within the generated range.
\end{proof}

\begin{theorem}[Patience Sorting]
The LNDS of a sequence of length $M$ can be computed in $O(M \log M)$ time by maintaining an array $\mathrm{tails}[]$ where $\mathrm{tails}[j]$ stores the smallest ending element of any non-decreasing subsequence of length $j+1$.
\end{theorem}

\begin{proof}
For each new element $v$, use \texttt{upper\_bound} to find the first position $p$ with $\mathrm{tails}[p] > v$. If $p$ is past the end, append $v$; otherwise set $\mathrm{tails}[p] = v$. The array $\mathrm{tails}[]$ remains non-decreasing throughout, each insertion costs $O(\log M)$ via binary search, and the final length of $\mathrm{tails}[]$ is the LNDS length by the standard theory of patience sorting.
\end{proof}

\section{Algorithm}

\begin{algorithm}[H]
\caption{Compute $S(n)$}
\begin{algorithmic}[1]
\State Generate stations: $(2^i \bmod n, 3^i \bmod n)$ for $i = 0,\dots,2n$
\State Remove duplicate stations
\State Sort stations lexicographically by $(x, y)$
\State Extract $y$-values into array $Y$
\State $\mathrm{tails} \gets []$
\For{each $v$ in $Y$}
    \State $p \gets \textsc{UpperBound}(\mathrm{tails}, v)$
    \If{$p = |\mathrm{tails}|$}
        \State Append $v$ to $\mathrm{tails}$
    \Else
        \State $\mathrm{tails}[p] \gets v$
    \EndIf
\EndFor
\State \Return $|\mathrm{tails}|$
\end{algorithmic}
\end{algorithm}

\begin{algorithm}[H]
\caption{Main: Compute $\sum_{k=1}^{30} S(k^5)$}
\begin{algorithmic}[1]
\State $\text{total} \gets 0$
\For{$k = 1$ to $30$}
    \State $\text{total} \gets \text{total} + S(k^5)$
\EndFor
\State \Return total
\end{algorithmic}
\end{algorithm}

\section{Complexity Analysis}

\begin{itemize}
  \item \textbf{Time:} $O(M \log M)$ per value of $n$, where $M \le 2n+1$. Total: $O\!\left(\sum_{k=1}^{30} k^5 \log k\right)$, dominated by $k = 30$ with $n = 24\,300\,000$.
  \item \textbf{Space:} $O(M)$ for storing stations and the $\mathrm{tails}$ array.
\end{itemize}

\section{Answer}

\[
\boxed{9936352}
\]

\end{document}
